\section{Introduction}

Social agents rarely act on a single, stable model of the social world; instead, they maintain different expectations for different partners. One partner may be familiar and reliable, another newly encountered and uncertain, and another predictable but adversarial. Repeated social exchange depends on learning such partner-specific regularities, including reciprocity, reputation, moral character, and social value~\cite{Behrens2008SocialLearning,Delgado2005MoralCharacter,KingCasas2005Trust}. In such settings, the question is not only what an agent believes about each partner, but how confidently those beliefs should guide policy inference.

Computational models of social decision-making can therefore be separated into three problems. The first is inferring a partner's behavioral disposition: whether their actions are cooperative, exploitative, reciprocal, or volatile. The second is modeling what they know or intend: the theory-of-mind problem of how another agent represents the world and plans action~\cite{FrithFrith2012,Hampton2008Mentalizing,Pitliya2025ToMActiveInference}. The third is estimating whether the resulting partner model is reliable enough to guide confident policy commitment~\cite{Schoeller2021TrustExtendedControl}. To address this third problem, an agent must decide not only what it believes about a partner, but how strongly to act on those beliefs. However, many formal models specify the content of social inference while treating confidence in those inferences as fixed, global, or implicit.

The active inference framework~\cite{Parr2022ActiveInference} provides a natural formalism for this problem because it links perception, belief updating, and policy selection through a common inferential process. In active inference, agents infer hidden states of the world and select policies by minimizing expected free energy. Moreover, policy inference is precision-weighted: policy precision controls how sharply an agent commits to its inferred policy distribution, making active inference a model not only of what an agent believes but how confidently those beliefs are expressed in policy~\cite{Friston2017ProcessTheory}.

Existing social active-inference models address complementary parts of the problem. Theory-of-mind models address how agents model what others know or intend through recursive belief structures and factorized generative models~\cite{Pitliya2025ToMActiveInference,RuizSerra2025Factorised}. Volatility and social-learning models address how prediction errors should change belief updating under uncertainty~\cite{Behrens2008SocialLearning,Diaconescu2017Hierarchical,Mathys2011HGF}. Social-learning work also shows that agents revise expectations from discrepancies between expected and observed behavior of others~\cite{joiner2017social}. Formal models of interpersonal inference have similarly applied active inference to repeated trust-game exchange, in which agents update beliefs about social partners and select policies under those beliefs~\cite{Moutoussis2014InterpersonalInference}. Related opponent-modeling and type-based approaches maintain hypotheses about other agents' behavioral strategies, goals, or beliefs in order to predict their behavior~\cite{AlbrechtStone2018ModellingAgents}. These approaches clarify how agents infer partners, model others' mental states, and revise beliefs under uncertainty. However, they do not by themselves explain how confidence in a social model should remain local to the relationship that generated the evidence and modulate policy commitment accordingly.

Affect provides a natural candidate mechanism for this confidence-calibration role. In psychology and neuroscience, affect is often described in terms of valence and arousal---the felt pull toward or away from situations and social commitments~\cite{Barrett2006EmotionParadox}. Free-energy and active-inference accounts extend this view by linking affective dynamics to prediction error, model evidence, emotional-state inference, and precision-weighted policy inference~\cite{Hesp2021DeeplyFeltAffect,joffily2013emotional,pattisapu2024free,smith2019simulating}. Prior affective-precision accounts make confidence endogenous by relating affect to subjective model fitness and policy precision~\cite{Hesp2021DeeplyFeltAffect}. On this view, affect can regulate how strongly current beliefs are expressed in policy, shaping the degree of commitment rather than the content of inference. In this work, \emph{affective precision} refers specifically to confidence in deploying a partner model during policy selection. It concerns how strongly partner beliefs guide policy rather than felt valence, emotion concepts, attachment, appraisal, or affective experience more broadly.

Related hierarchical active-inference work has also modeled metacognitive states that regulate precision within the generative hierarchy~\cite{SandvedSmith2021MentalAction}. Here, we focus specifically on making such confidence relationship-specific and linking it to the predictive reliability of individual partner models.

We therefore define partner-local affective precision as a relationship-specific confidence signal. For each partner \(k\), the focal agent maintains a posterior \(q(\beta_k)\) over an inverse policy-precision variable \(\beta_k\), such that lower \(\beta_k\) corresponds to greater confidence and sharper policy selection. Evidence from partner \(k\) updates only that partner's confidence estimate, and its posterior mean modulates policy precision for policies evaluated under the corresponding partner model. Unlike prior affective-precision accounts centered on policy-inference quantities~\cite{Hesp2021DeeplyFeltAffect}, we derive the affective signal from the predictive evidence assigned to the partner's observed response under the focal agent's current model. The signal therefore tracks whether the current model predicts this partner's behavior rather than whether the interaction produces a favorable payoff, preserving relationship-specific differences that a global confidence estimate would collapse.

We evaluate this mechanism in a repeated multi-partner graded trust game with partner-choice conditions, abrupt betrayal, and profile variation. The remainder of this paper is organized as follows: Section~\ref{sec:methods} specifies the multi-partner trust game, per-partner generative models, and the partner-local affective precision mechanism. Section~\ref{sec:results} reports simulation results and analyses. Section~\ref{sec:discussion} discusses implications, limitations, and directions for future work.
